\section{Фильтрация изображений с периодичностью}
Скачаем изображение 1 с \href{https://drive.google.com/drive/folders/1oewu85taKvyxAUhXNH48ImTFtI1SpFag}{этого гугл-диска} (см. рис. \ref{1}).
Видим  периодичность с углом около $40^\circ$.
\begin{figure}[H]
    \centering
    \includegraphics[width=0.55\linewidth]{/home/ns/Desktop/ITMO-courses/Fourier_Methods/Lab6/1.png}
    \caption{Картинка}
    \label{1}
\end{figure}
Выполним следующие шаги:

\begin{itemize}
    \item Загрузим изображение в \texttt{Python} при помощи команды \texttt{cv2.imread()}.
    \item Преобразуем полученный массив к вещественному типу и поделим все значения на 255 - максимальное значение яркости цвета.
    \item Найдём двумерный Фурье-образ массива и сдвинем его в центр с помощью команд \texttt{fft.fftshift(fft.fft2())} для каждого из цветовых каналов.
    \item Разделим полученный образ на массивы модулей и аргументов с помощью функций \texttt{np.abs} и \texttt{np.angle}.
    \item Для удобства работы, найдём логарифм от массива модулей и нормализуем его значения в диапазон от 0 до 1.
    Чтобы избежать неопределённости в логарифме, предварительно прибавим ко всем значениям 1.
    \item Сохраним полученный массив (нормализованный логарифм модуля Фурье-образа) как изображение командой \texttt{cv2.imwrite}.
    \item Проанализируем полученное на предыдущем шаге изображение. Найдём пики, соответствующие периодичности на исходной картинке.
    \item Отредактируем полученный Фурье-образ: сгладим все ненужные цветовые пики, отвечающие за гармоники, от которых хотим избавиться.
    \item Восстановим картинку из отредактированного образа, проделав обратные шаги.
\end{itemize}
\newpage
Изображение логарифма модуля образа исходной картинки представлено на рис. \ref{2}.

\begin{figure}[H]
    \centering
    \includegraphics[width=0.6\linewidth]{/home/ns/Desktop/ITMO-courses/Fourier_Methods/Lab6/plots/fourier_spectrum_1.png}
    \caption{Логарифм модуля образа}
    \label{2}
\end{figure}

Мы видим множество пиков, расположенных по диагонали на картинке. Выделим их, уберём и посмотрим, что получится (см. рис.\ref{3}).
\begin{figure}[H]
    \centering
    \includegraphics[width=0.45\textwidth]{/home/ns/Desktop/ITMO-courses/Fourier_Methods/Lab6/plots/fourier_spectrum_1_red.png}
    \hfill
    \includegraphics[width=0.45\textwidth]{/home/ns/Desktop/ITMO-courses/Fourier_Methods/Lab6/plots/fourier_spectrum_1_filtered.jpg}
    \caption{Выделенные области с пиками на спектре (слева) и 
    очищенный от пиков спектр (справа)}
    \label{3}
\end{figure}

Сравним теперь исходное и отфильтрованное изображение (см. рис. \ref{4}).\\
Видим, что периодичность исчезла почти полностью, но есть и побочный эффект - изображение стало более тусклым.
Это происходит потому что, затирая пики, мы задеваем часть остальной информации изображения, от чего оно искажается.
\begin{figure}[H]
    \centering
    \includegraphics[width=0.45\textwidth]{/home/ns/Desktop/ITMO-courses/Fourier_Methods/Lab6/1.png}
    \hfill
    \includegraphics[width=0.45\textwidth]{/home/ns/Desktop/ITMO-courses/Fourier_Methods/Lab6/plots/periodic_filtered_image.png}
    \caption{Исходное изображение (слева) и отфильтрованное (справа)}
    \label{4}
\end{figure}
\textbf{Выводы}\\
Используя многомерное преобразование Фурье, можно осуществлять фильтрацию на изображениях, в данном случае, мы видим,
как можно убирать гармоники из спектра, создающие периодические структуры на изображении. При этом чем сильнее мы
гасим пики на спектре, тем больше информации мы задеваем, и изображение немного искажается (меняются оттенки), поэтому нельзя
затирать их слишком сильно, нужно искать "золотую середину", когда периодичность убрана достаточно хорошо, и при этом нет 
сильного изменения цветов на изображении, чтобы оно не стало совсем тусклым.